\documentclass[11pt]{article}
\usepackage[a4paper,margin=0.85in]{geometry}
\usepackage{amsmath,amssymb,amsthm,mathtools}
\usepackage{graphicx}
\usepackage{float}
\usepackage[hidelinks]{hyperref}
\usepackage{microtype}

\newtheorem{theorem}{Theorem}
\newtheorem{corollary}[theorem]{Corollary}
\newtheorem{lemma}[theorem]{Lemma}
\theoremstyle{remark}
\newtheorem{remark}[theorem]{Remark}

\newcommand{\R}{\mathbb R}
\newcommand{\F}{\mathcal F}
\newcommand{\cE}{\mathcal E}

\title{A uniform \(L^1\) semigroup bound for homogeneous elliptic operators}
\author{Candidate full solution packet for Question 4.5 of arXiv:1805.08557}
\date{26 August 2026}

\begin{document}
\maketitle

\begin{center}
\fbox{\parbox{0.91\textwidth}{\textbf{Status.}
Candidate full solution, likely valid, awaiting specialist review.
In the constant-coefficient setting of the source paper, Question 4.5 has an
affirmative answer in the fastest listed decay regime, with \(\beta=0\).  More
precisely, the semigroup is uniformly bounded on \(L^1(\R^d)\).}}
\end{center}

\section{Source question and setting}

Ben-Artzi and Einav consider \(H=P(D)\) on \(L^2(\R^d)\), where in the
homogeneous elliptic subsection
\[
 P(\xi)=\sum_{|\alpha|=m}a_\alpha\xi^\alpha,
 \qquad m\in\{2,4,6,\ldots\},
\]
has real coefficients and satisfies their Assumption A3.  In particular,
\(P\) is homogeneous of degree \(m\), \(P(0)=0\), and
\[
 P(\xi)\ge c|\xi|^m\qquad(\xi\in\R^d)
\]
for some \(c>0\).  The semigroup generated by \(-P(D)\) is denoted here by
\(T_t=e^{-tP(D)}\); it is the source's \(P_t\).

Question 4.5 asks whether for every such operator and every relevant initial
datum \(u\) there are \(C_2=C_2(u)\ge0\) and \(\beta\in\R\) such that
\[
 \|T_tu\|_1^2\le \|u\|_1^2+C_2t^\beta
 \qquad(t\ge0) .
\]
An affirmative answer with \(C_2=0\) or \(\beta\le0\) gives the source's
fastest listed decay rate, \(O(t^{-d/m})\) \cite{BenArtziEinav}.

\begin{figure}[H]
 \centering
 \includegraphics[width=0.98\textwidth]{figures/open_problem_crop.png}
 \caption{Question 4.5 and its stated consequence on page 14 of the source
 paper.  The crop preserves the full rendered page width.}
\end{figure}

\section{Proof intuition}

Homogeneity makes every positive-time heat kernel a dilation of a single
kernel.  Indeed, if
\(k=\F^{-1}(e^{-P})\), then the kernel at time \(t\) is
\[
 k_t(x)=t^{-d/m}k(t^{-1/m}x).
\]
Ellipticity forces \(e^{-P}\) to decay exponentially together with all of its
derivatives, so \(k\) is a Schwartz function.  Thus \(k\in L^1\), while the
dilation formula gives \(\|k_t\|_1=\|k\|_1\) for every \(t>0\).  Young's
inequality then yields a time-uniform \(L^1\) operator norm.  No positivity of
the higher-order heat kernel is required.

\section{The uniform kernel estimate}

\begin{lemma}\label{lem:kernel}
Let \(P\) be a real homogeneous polynomial of even degree \(m\) satisfying
\(P(\xi)\ge c|\xi|^m\).  Set
\[
 k=\F^{-1}(e^{-P}),\qquad
 k_t=\F^{-1}(e^{-tP})\quad(t>0).
\]
Then \(k\in\mathcal S(\R^d)\),
\[
 k_t(x)=t^{-d/m}k(t^{-1/m}x),
 \qquad
 \|k_t\|_{L^1}=\|k\|_{L^1}.
\]
\end{lemma}

\begin{proof}
For every multi-index \(\gamma\), repeated differentiation writes
\(\partial^\gamma(e^{-P})\) as a polynomial times \(e^{-P}\).  Since
\(P(\xi)\ge c|\xi|^m\), multiplication by any further monomial still leaves
a bounded function.  Hence \(e^{-P}\in\mathcal S(\R^d)\).  The Fourier
transform is an automorphism of the Schwartz space, so
\(k\in\mathcal S(\R^d)\subset L^1(\R^d)\).

Homogeneity gives
\[
 e^{-tP(\xi)}=e^{-P(t^{1/m}\xi)}.
\]
The Fourier scaling rule therefore yields
\(k_t(x)=t^{-d/m}k(t^{-1/m}x)\).  A change of variables gives the asserted
equality of \(L^1\)-norms.
\end{proof}

\section{Full answer to Question 4.5}

\begin{theorem}\label{thm:answer}
In the constant-coefficient homogeneous elliptic setting of Question 4.5,
the semigroup \(T_t=e^{-tP(D)}\) extends to a uniformly bounded convolution
semigroup on \(L^1(\R^d)\).  Specifically, with
\[
 M_P:=\bigl\|\F^{-1}(e^{-P})\bigr\|_{L^1}<\infty,
\]
one has
\[
 \boxed{\ \|T_tu\|_{L^1}\le M_P\|u\|_{L^1}
 \quad\text{for all }u\in L^1(\R^d),\ t>0.\ }
\]
Consequently Question 4.5 has an affirmative answer with
\[
 \beta=0,
 \qquad
 C_2(u)=\max\{M_P^2-1,0\}\,\|u\|_{L^1}^2.
\]
\end{theorem}

\begin{proof}
On Schwartz functions, Fourier diagonalization of the constant-coefficient
operator gives
\[
 T_tu=\F^{-1}(e^{-tP}\widehat u)=k_t*u.
\]
By Lemma~\ref{lem:kernel} and Young's convolution inequality,
\[
 \|T_tu\|_1\le\|k_t\|_1\|u\|_1=M_P\|u\|_1.
\]
This estimate extends the formula uniquely from Schwartz functions to all of
\(L^1\), and on \(L^1\cap L^2\) it agrees with the original \(L^2\)
spectral semigroup.  The multiplier identity gives the semigroup law first on
Schwartz functions and then by density on \(L^1\).  Moreover,
\(\int k=1\), and the displayed dilations of \(k\) form a bounded approximate
identity, so the extension is strongly continuous at \(t=0\).

Squaring the estimate and using the displayed choice of \(C_2(u)\) proves
\[
 \|T_tu\|_1^2\le \|u\|_1^2+C_2(u)t^0\qquad(t>0).
\]
At \(t=0\), \(T_0u=u\), so the source inequality also holds there (whether
one defines \(0^0=1\) or simply checks the endpoint separately).
\end{proof}

\begin{remark}
If the semigroup is positivity preserving and mass preserving, then its
convolution kernel has \(L^1\)-norm one and the stronger contraction estimate
\(\|T_tu\|_1\le\|u\|_1\) holds.  The proof above covers higher-order kernels
that change sign, with a fixed constant \(M_P\).
\end{remark}

\section{The resulting decay rate}

The source proves the homogeneous Nash-type inequality
\[
 \|v\|_2^2
 \le C_P\,\cE(v)^{d/(m+d)}\|v\|_1^{2m/(m+d)},
 \qquad
 \cE(v)=\|P(D)^{1/2}v\|_2^2.
\]
For completeness, the uniform \(L^1\) estimate also gives the stated decay
directly.  For \(u\in L^1\cap L^2\), put \(y(t)=\|T_tu\|_2^2\).  The spectral
calculus gives \(y'(t)=-2\cE(T_tu)\).  Rearranging the Nash inequality and
using Theorem~\ref{thm:answer} yields
\[
 y'(t)
 \le -c_P\,(M_P\|u\|_1)^{-2m/d}y(t)^{1+m/d}.
\]
Integration gives
\[
 y(t)\le
 \left(y(0)^{-m/d}
 +c'_P(M_P\|u\|_1)^{-2m/d}t\right)^{-d/m}
 =O(t^{-d/m}).
\]
Thus the answer realizes the \(\beta\le0\) line predicted immediately after
Question 4.5.

\section{Verification, scope, and novelty check}

The proof uses only four auditable facts: the source's operator is a Fourier
multiplier; ellipticity makes \(e^{-P}\) Schwartz; homogeneity gives exact
kernel scaling; and Young's inequality converts this into an \(L^1\) bound.
The quantifiers are stronger than requested: \(M_P\) is independent of both
\(t\) and \(u\), while the allowed additive constant is then chosen for each
\(u\).  The argument does not assume positivity, and it handles \(t=0\)
separately.

The conclusion is scoped to the source's Section 4 setting: scalar,
constant-coefficient operators \(P(D)\) on \(\R^d\), with a real positive
homogeneous polynomial symbol.  It does not claim the same bound for
variable-coefficient elliptic operators or for arbitrary nonhomogeneous
symbols.

A bounded literature check through 26 August 2026 searched the run's
registry, solution, attempt, and proof-gap indexes; arXiv:1805.08557 and its
exact Question 4.5 wording; and primary-source searches combining
``homogeneous elliptic operator'', ``heat kernel'', ``\(L^1\)'', and
``semigroup''.  It found the source paper and general heat-kernel literature,
but no later paper explicitly identifying itself as an answer to Question
4.5.  The kernel-scaling ingredients are classical, so this search does not
establish historical priority.  The packet's claim is the direct
identification and complete proof of the answer to the source question.

The highest-value specialist check is the interpretation of ``homogeneous
elliptic operator'' exactly as fixed by the surrounding constant-coefficient
subsection.  Under that stated context, the proof is complete.

\begingroup
\small
\begin{thebibliography}{9}
\bibitem{BenArtziEinav}
J. Ben-Artzi and A. Einav,
\emph{Weak Poincar\'e inequalities in the absence of spectral gaps},
Annales Henri Poincar\'e 21 (2020), 359--375;
arXiv:1805.08557v5.

\bibitem{Grafakos}
L. Grafakos,
\emph{Classical Fourier Analysis}, 3rd ed.,
Graduate Texts in Mathematics 249, Springer, 2014.
\end{thebibliography}
\endgroup

\end{document}
